\def\anon{1} %% set to 1 for anonymous submissions, hides acknowledgements and author names
\def\full{1} %% set to 0 for springer proceedings


\ifnum\full=1
\documentclass[11pt]{llncs}


\addtolength{\parskip}{1pt}
\else
\documentclass[10pt, runningheads]{llncs}
\usepackage{times}
\fi


\input{../header.tex}


\begin{document}


%\title{Notes: On the Resilience Order of Weightwise Almost Perfectly Balanced Functions}
\title{On the Resilience Order of W(A)PB Functions}


%\subtitle{}



	\titlerunning{On the Resilience Order of W(A)PB Functions}
	\author{
			\mbox{Martin Grenouilloux \orcidID{0009-0003-1998-6092}},
			\mbox{Chunlei Li \orcidID{0000-0002-3792-769X}},
		\mbox{Pierrick M\'eaux\orcidID{0000-0001-5733-4341}}
	}

	\authorrunning{ M. Grenouilloux, C. Li, P. M\'eaux}

	\institute{University of Bergen, Bergen, Norway\\
	\email{martin.grenouilloux@uib.no,Chunlei.Li@uib.no}\\
University of Luxembourg, Luxembourg\\	\email{pierrick.meaux@uni.lu}
}






	%----------------------------------------------------------------
	\maketitle



 \setcounter{page}{1}

\begin{abstract}
The recent development of Fully Homomorphic Encryption (FHE)
witnessed the emergence of a new generation of tailored cryptographic primitives
designed to meet its specific criteria. Among promising candidates for FHE
constructions stands out the FLIP cipher, which employs Boolean functions that
are evaluated only on specific subsets of $\mathbb{F}_2^n$.
In this article, we study Weightwise (Almost) Perfectly Balanced (W(A)PB) Boolean functions, which are balanced on each of these subsets. While W(A)PB functions have been of great
interest for new constructions recently, some of their aspects like resilience
remain uncertain. As such, we take a first step at characterizing the resilience
of W(A)PB functions, through their properties as correctors. We highlight its
close connection with the (restricted) Walsh transform and uncover an algebraic
relation between Krawtchouk matrices and Vandermonde matrices, which
reduces the problem of determining the correction order of a W(A)PB function to a weak
instance of the Prouhet-Tarry-Escott problem.
This reduction helps us show that
for infinitely many integers $n$, W(A)PB functions in $n$ variables can have correction order at most the Hamming weight of $n$ minus one.
We make a conjecture that this observation holds for any positive integer $n$, which is verified for $n$ up to $62$.

\end{abstract}

\keywords{WAPB functions, resilience order, correction order, Krawtchouk polynomials}

\end{document}
